\documentclass[./dissertation.tex]{subfiles}
\begin{document}

\chapter{Experiments}
\label{Experiments}

To evaluate the proposed method, experiments were done on a sorting task for recyclable items with a robotic manipulator. The following experiments are the main contributions of this paper. 
% Add something about metrics evaluated 

\section{Task Description}
The proposed task is to autonomously sort recyclable items within a constrained workspace. The item set consists of aluminum cans and plastic bottles with varying degrees of deformation, which introduces variability in shape and grasp difficulty. The object manipulation is performed with a KUKA LBR iisy 3 r760 collaborative robotic arm equipped with an Ubiros Gentle Pro gripper at the end effector. The context of the workspace is captured using an overhead Oak-D S2 RGB-D camera, which provides synchronized color and depth images of the scene. These images are processed and encoded to form the input to the learned policy. The policy outputs continuous control actions for the manipulator, specifying the relative end effector placement. After executing the policy action with the corresponding motion trajectory, the robot automatically closes the gripper to secure the object, moves to the recycling bin, and opens the gripper to release any grasped object. An action is considered successful if the robot is able to securely grasp the assigned object and transport it to the recycling bin. The experimental setup is illustrated in Figure \ref{Workspace} and the overall sorting procedure is summarized in Algorithm \ref{Sorting Loop}.

\begin{figure}[h]
\centering\includegraphics[width=0.375\textwidth]{figures/Workspace Labeled.png}
\caption{Experimental Setup. A KUKA robotic manipulator is used to grab aluminum cans and plastic bottles located on a conveyor belt. An Oak-D RGB-D camera captures images of the scene, from which the 3D position and orientation of each item is calculated. With this information, the manipulator executes grasping actions to move the recyclable items to a recycling bin}
\label{Workspace}
\end{figure}

\begin{algorithm}
\caption{Sorting}
\label{Sorting Loop}
\textbf{Input}: Learned policy $\pi_\phi$
\begin{algorithmic}[1]
\While {items in workspace}
    \State Camera captures RGB and Depth images of the workspace
    \State Process images to obtain context $x$ and encoded context $\hat{x}$
    \State Robot moves to "Home" position
    \State Determine action from policy: $a \gets \pi_\phi(\hat{x})$
    \State Robot executes action $a$ by moving to relative end effector position
    \State Robot closes gripper to grab item
    \State Robot moves to "Bin" position
    \State Robot opens gripper to drop item
\EndWhile
\end{algorithmic}
\end{algorithm}

\section{Context Estimation}
The context of the workspace is obtained using images from an overhead RGB-D camera. When the current context is requested, the camera captures an RGB image and a corresponding depth image of the workspace. The RGB image is passed through a YOLO v11 instance segmentation model, which produces pixel-wise masks of the aluminum cans and plastic bottles present in the scene. The YOLO model was trained on a combination of publicly available datasets containing annotated images of general cans and bottle, as well as a custom dataset containing images collected directly from the workspace with manual annotation. Incorporating general and task-specific data enables the model to generalize and maintain accuracy for specific objects with the instance segmentation procedure. 

After obtaining the segmented instances from the RGB image, Principal Component Analysis (PCA) \parencite{PCA} is performed on the mask of each item. PCA provides an efficient method to calculate different geometric properties from the object masks, and is first utilized to compute the centroid. The centroid represents the average position of all the pixels in the object mask, effectively providing a 2D estimate of the object's center in the image plane. In addition to the centroid, PCA is utilized to calculate the principal axis of each mask, which is used to indicate the orientation of the each object in the image. Figure \ref{State Estimation} provides a visual overview of this process, showing the initial RGB image, the extracted object masks, and the result of the PCA analysis on each mask.

The corresponding depth image is then used to recover the 3D position of the centroid from its 2D location in the image. Using the camera intrinsics and a standard pinhole model, the pixel corresponding to the centroid can be projected along a 3D ray. The depth value at the centroid's pixel represents the distance along this ray, allowing the computation of the object's centroid relative to the camera. This 3D position is transformed from the camera's local coordinate frame to the robot's global coordinate frame using ROS2 transformations. 

The principal axis obtained from PCA defines the object's yaw angle in the robot frame. However, for a complete 6D pose representation, the pitch angle must also be estimated. To compute the pitch, a two-point method is employed. Two points are selected near opposite ends of the object along the principal axis. The 3D positions of these points are recovered from the depth image using the same strategy applied to the centroid. Let the vertical difference between the two points be $\Delta z$, and let the planar displacement in the $xy$-plane be $\Delta xy$. The pitch angle is then computed by $\vartheta=\arctan{\frac{\Delta z}{\Delta xy}}$, which provides a measure of the object's incline relative to the workspace plane. Due to assumed axial symmetry of the bottles and cans, the roll angle is assumed to be zero. Therefore, the full 6D pose of the each object is defined by the 3D position of the centroid and the estimated pitch and yaw angles.

\begin{figure}[h]
    \centering
    \begin{minipage}[c]{0.32\textwidth}
        \centering
        \includegraphics[width=\textwidth]{figures/OG Bottles and Cans.png}
        \subcaption{}
        \label{OG}
    \end{minipage}
    \hspace{0.01\textwidth}
    \begin{minipage}[c]{0.32\textwidth}
        \centering
        \includegraphics[width=\textwidth]{figures/Masks.png}
        \subcaption{}
        \label{Masks}
    \end{minipage}
    \hspace{0.01\textwidth}
    \begin{minipage}[c]{0.32\textwidth}
        \centering
        \includegraphics[width=\textwidth]{figures/PCA.png}
        \subcaption{}
        \label{PCA}
    \end{minipage}
    \caption{Context Estimation. Figure~\ref{OG} shows the original RGB image. This image is passed through a YOLO v$11$ instance segmentation model to obtain masks of the bottles and cans, as shown in Figure~\ref{Masks}. Principal Componet Analysis (PCA) is performed on the masks to determine the centroid and principal axes, shown in Figure~\ref{PCA}}
    \label{State Estimation}
\end{figure}

\section{Context Definition}
Because the robot focuses on grasping one item at a time, the context representation only includes the 6D pose of the selected target. This approach greatly reduces the context space by concentrating on the object that is relevant for the current grasping action, rather than including the positions and orientations of all other objects, which is difficult when the number of items is dynamic. Therefore, the context captures the position and orientation of the active object in 3D space, along with basic contextual information about the scene. Formally the context is defined as:

\[
\mathbf{x}=[x , y, z, \omega, \psi, t, n],
\]

where $x,y,z \in \mathbb{R}$ represent the Eucledian coordinates of the object's centroid, $\omega$ and $\psi$ denote the pitch and yaw angles describing the object's orientation, $t$ represents the object type (e.g., can or bottle), and $n$ denotes the number of objects currently present in the scene. Information about the object type $t$ is necessary as the geometry of aluminum cans and plastic bottles differs. 

Due to the physical constraints of the workspace, the following bounds are imposed on the context variables: 

\[
\begin{gathered}
x \in [0.20, 0.55], \
y \in [-0.35, 0.35], \
z \in [0.80, 0.95],  \\
\omega \in [-\frac{\pi}{2}, 0], \
\psi \in [-\pi, \pi].
\end{gathered}
\]

These bounds ensure that the context representation includes only feasible poses, preventing the robot from considering impossible positions or orientations during planning or learning.

Because the true context space is continuous, there are infinitely many possible contexts. To significantly reduce the size of the context space for the reward and policy networks, a discretized encoding is introduced. The workspace is first divided into two equal partitions along the $x$ and $y$ axes, producing four spatial quadrants. Each object is assigned a quadrant based on the position of its centroid, compressing all possible continuous $x$ and $y$ values into four discrete values. Then, the pitch angle $\omega$ range is divided into three uniform bins, and the yaw angle $\psi$ range is divided into four uniform bins. Like the quadrant representation, the pitch and yaw bins compress the context representation into a coarse but still meaningful form. The object type $t$ is preserved in the discretized context to maintain the distinction between different recyclable items, while the number of objects $n$ is removed to simplify the encoding and because only a single target is considered.

The resulting discretized state is compactly defined as:

\[
\mathbf{\hat{x}} = [q, \Omega, \Psi, t],
\]

where $q \in \{ 1,2,3,4 \}$ denotes the spatial quadrant of the object, $\Omega \in \{ 1,2,3 \}$ represents the discretized pitch bin, $\Psi \in \{ 1,2,3,4 \}$ denotes the discretized yaw bin, and $t$ represents the object type. This encoding provides a simplified yet informative summary of the object and reduces the continuous context to a compact representation that is suitable for reward modeling and policy learning. Figure \ref{Discretized Workspace} shows the discretized workspace along with an example of an item mapped to the corresponding quadrant and orientation bins.

\begin{figure}[h]
    \centering
    \begin{minipage}[c]{0.49\textwidth}
        \centering
        \includegraphics[width=\textwidth]{figures/Quadrants.png}
        \subcaption{}
        \label{Quadrants}
    \end{minipage}
    \hspace{0.01\textwidth}
    \begin{minipage}[c]{0.49\textwidth}
        \centering
        \includegraphics[width=\textwidth]{figures/Slices.png}
        \subcaption{}
        \label{Yaw Bin}
    \end{minipage}
    \\[0.25cm]
    \begin{minipage}[c]{0.5\textwidth}
        \centering
        \includegraphics[width=\textwidth]{figures/Sectors.png}
        \subcaption{}
        \label{Pitch Bin}
    \end{minipage}
    \caption{Discretized Workspace: Figure~\ref{Quadrants} shows the xy-plane of the workspace divided into four quadrants. In this example, the bottle is located in the second quadrant. Figure~\ref{Yaw Bin} illustrates the discretization of the possible yaw angles into four bins; the bottle falls within the first bin in this example. Similarly, the possible pitch angles are discretized into three bins, with the bottle located in the first bin, as shown in Figure~\ref{Pitch Bin}.}
    \label{Discretized Workspace}
\end{figure}

\section{Action Definition}
The actions for the sorting task define how the robotic manipulator moves prior to attempting a grasp. Each action specifies a desired relative displacement of the tool center point (TCP), i.e., the center of the gripper, with respect to the current object pose. To execute an action, the robot computes an inverse kinematics \parencite{IK} solution and plans a motion to the specified pose. The action is 6-dimensional, continuous, and is defined by the following 

\[
a = [\Delta l, \Delta w, \Delta h, \Delta \varphi, \Delta \omega, \Delta \psi]
\]

where $\Delta l, \Delta w, \Delta h$ are the relative Cartesian offsets in meters along the object's local XYZ frame, and $\Delta \varphi, \Delta \omega, \Delta \psi$ represent the relative orientation offsets (roll, pitch, and yaw, respectively) in radians. All position and orientation components are defined relative to the pose of the target object. This means that the action does not command an absolute pose in the workspace frame, but rather a pose expressed in the object's local frame. For example, the action $[0.0, 0.0, 0.0, \frac{\pi}{4}, 0.0, 0.0]$ moves the TCP to the object's centroid with matching pitch and yaw, while introducing a relative roll offset of $45^\circ$. Each component of the relative action is clamped to pre-defined bounds defined by

\begin{figure}[h]
\centering\includegraphics[width=0.5\textwidth]{figures/Local Axis.png}
\caption{Local Coordinate Frame. The yellow dot on the bottle represents its centroid, while the red, green, and blue axes indicate the positive directions of the object's local $x$, $y$, and $z$ axes, respectively. An action containing a $+\Delta l$ component would produce a translation along the red axis, starting from the yellow dot.}
\label{Local Axis}
\end{figure}

\[
\begin{gathered}
\Delta l \in [-0.10, 0.10], \
\Delta w \in [-0.10, 0.10], \
\Delta h \in [-0.01, 0.10],  \\
\Delta \varphi \in [-\frac{\pi}{4}, \frac{\pi}{4}], \
\Delta \vartheta \in [-\frac{\pi}{4}, \frac{\pi}{4}], \
\Delta \psi \in [-\frac{\pi}{4}, \frac{\pi}{4}]
\end{gathered}
\]

to ensure the manipulator only attempts feasible grasps.

Although the policy outputs a relative action, execution requires an absolute pose in the workspace frame. The absolute action is computed with component-wise addition between the action and the true state defined in the previous section. Before execution, the resulting absolute target pose is clamped to the same bounds defined for the workspace state.

\section{Reward Definition}
Rewards are assigned autonomously to each action executed in the workspace. The reward is binary: actions that result in a successful grasp receive a reward of 1, while unsuccessful actions receive a reward of 0. Action success is determined using two criteria: a reduction in the total workspace items, and the measured current load of the gripper. If the number of items in the workspace decreases by one following the action and gripper current exceeds 150 mA, the action is classified as successful and assigned a reward of 1. Otherwise, the action is classified as a failure and assigned a reward of 0. This dual-condition check ensures that grasp attempts that fail to grab the object are not incorrectly labeled as successful.

To train the reward model, Binary Cross-Entropy (BCE) loss is used. Cross-Entropy loss \parencite{20} is widely used for probability models and classification tasks. Because the reward data takes on one of two labels (0 or 1), the binary variant of Cross-Entropy is appropriate. The Binary Cross-Entropy is defined as

\[
\mathcal{L}_{BCE}(\theta) =
-\frac{1}{B}\sum_{b=1}^B
y_b\cdot \log(\hat y_b) + 
(1-y_b) \cdot \log(1-\hat y_b)
\]

where $y$ denotes the true labels of input $x$ and $\hat y$ represents the predicted labels  produced by a function $f_\theta(x)$. Minimizing this loss encourages the reward model to assign high probability to successful actions and low probability to unsuccessful ones. 

A binary reward setup is effective for evaluating the method because multiple actions can yield successful grasps and achieve the maximum reward. However, many of these actions are suboptimal due to inherent uncertainty in the workspace and the grasping process. Rather than optimizing for any action that simply achieves success, the modified objective encourages convergence towards actions that are more optimal and robust to uncertainty. These optimal actions are often located near the center of regions with consistently high rewards. Converging to these regions improves overall performance and helps avoid policies that settle in local optima or rely on highly situational actions.

\section{Initial Dataset}
As mentioned earlier, the goal of the initial data collection phase is to produce a dataset with diverse states, actions, and rewards. To achieve this, data was collected from both an expert and a random policy to ensure a mix of successful and unsuccessful interactions. Expert data was generated by hand-guiding the KUKA collaborative robot to a variety of feasible grasp poses for the assigned object. With this teaching process, 160 expert demonstrations were produced, corresponding to five demonstrations per state. Additionally, a random policy was implemented, sampling actions uniformly within the action bounds. Executing this policy yielded 340 additional interactions with mainly unsuccessful attempts to provide negative reward examples. Together, the expert and random data form an initial dataset of 500 interactions, containing both high quality grasps and exploration failures to produce an effective initial reward model.

\end{document}
